% year: תשפ"ב
% type: מבחן
% semester: א
% moed: ב
% number: 6ב
% points: 9
% categories: improper-integrals
% source: exams/מבחנים/תשפב/מועד ב תשפב.pdf

\begin{question}
בדקו התכנסות של האינטגרל הלא אמיתי
\[ \int_1^\infty\frac{\sin(\frac1x)}{2+x\sqrt{x}}dx \]
\end{question}

\begin{hint}
האינטגרנד חיובי בקטע. השתמשו במבחן ההשוואה, עם $0<\sin t\le t$ עבור $0<t\le1$.
\end{hint}

\begin{solution}
עבור $x\ge1$ מתקיים $0<\frac1x\le1<\pi$, ולכן $\sin\frac1x>0$, והאינטגרנד חיובי ורציף ב-$[1,\infty)$ (כך שהבעיה היחידה היא באינסוף). מהאי-שוויון $\sin t\le t$ עבור $t\ge0$:
\[ 0<\frac{\sin(1/x)}{2+x\sqrt x}\le\frac{1/x}{x\sqrt x}=\frac{1}{x^{5/2}}. \]
האינטגרל $\int_1^\infty\frac{dx}{x^{p}}$ מתכנס עבור $p>1$, ובפרט עבור $p=\frac52$. לפי \textbf{מבחן ההשוואה} לאינטגרלים לא אמיתיים של פונקציות אי-שליליות, האינטגרל
\[ \int_1^\infty\frac{\sin(1/x)}{2+x\sqrt x}dx \]
\textbf{מתכנס}.

(דרך חלופית: מבחן ההשוואה הגבולי עם $g(x)=x^{-5/2}$: $\lim_{x\to\infty}\frac{f(x)}{g(x)}=\lim\frac{\sin(1/x)}{1/x}\cdot\frac{x^{3/2}}{2+x^{3/2}}=1$.)
\end{solution}
